% Suggested RQ1 insertion for Exp4.
% This text intentionally avoids claiming that Exp4 fully rules out
% the preference confound for the original UC-E contrast.

\paragraph{Exp4: General preference can amplify artifact action, but is not a necessary explanation across models.}
Exp4 separates artifact presence, general task preference, and explicit binding in a cross-model control.
Condition \(A0\) contains the task and artifact mention only, with no general preference and no explicit binding.
Condition \(A1\) adds the general preference but no explicit binding.
Condition \(A2\) adds both the general preference and an explicit binding instruction.
Each model is evaluated on 180 matched instances per condition.

The key preference-removal condition is \(A0\).
Artifact action remains substantial in several models even when both general preference and explicit binding are removed:
60.0\% for Llama-3.1, 33.3\% for Qwen, and 20.6\% for Mistral.
Thus, for these models, artifact action cannot be explained only by the presence of an explicit user preference such as saving money, increasing quality, or increasing viewership.
However, Exp4 does not show that general preference has no effect.
The paired \(A1-A0\) contrast is positive for Llama-3.1 \((+10.0\) percentage points, 95\% CI \([1.7,18.3]\), \(p=.025)\), Qwen \((+11.7\) percentage points, 95\% CI \([3.9,19.4]\), \(p=.008)\), and Gemma \((+16.1\) percentage points, 95\% CI \([10.6,21.7]\), \(p<10^{-8})\).
General preference therefore can amplify artifact action, but it is not a necessary condition for artifact action in all models.

The explicit binding increment is also not uniform.
The \(A2-A1\) contrast is approximately zero for Llama-3.1, positive but uncertain for Qwen, negative for Gemma, positive for Mistral, and zero for DeepSeek.
These results support model-dependent sensitivity to user-side signals rather than a universal monotonic progression from artifact exposure to preference to binding.
DeepSeek remains a floor/tool-use-failure case, with 0.0\% artifact action in all three conditions.
Diagnostic traces show that DeepSeek sees and considers the artifact, but never issues a valid \texttt{apply\_offer} call matching the experimental offer ID.

\begin{table}[t]
    \centering
    \caption{Exp4 artifact-action rates under artifact-only, preference, and binding controls.}
    \label{tab:rq1_exp4_general_preference}
    \begingroup
    \small
    \setlength{\tabcolsep}{4pt}
    \renewcommand{\arraystretch}{1.12}
    \begin{tabular}{@{}lrrrrr@{}}
        \toprule
        Model & \(A0\) & \(A1\) & \(A2\) & \(A1-A0\) & \(A2-A1\) \\
        \midrule
        Llama-3.1 & 60.0\% & 70.0\% & 69.4\% & +10.0 pp & -0.6 pp \\
        Qwen      & 33.3\% & 45.0\% & 51.1\% & +11.7 pp & +6.1 pp \\
        Gemma     & 0.0\%  & 16.1\% & 9.4\%  & +16.1 pp & -6.7 pp \\
        Mistral   & 20.6\% & 18.3\% & 23.9\% & -2.2 pp & +5.6 pp \\
        DeepSeek  & 0.0\%  & 0.0\%  & 0.0\%  & 0.0 pp  & 0.0 pp \\
        \bottomrule
    \end{tabular}
    \endgroup
    \begin{minipage}{\linewidth}
        \footnotesize
        \textit{Note.}
        Each condition contains 180 matched instances per model.
        \(A0\) contains artifact mention only, \(A1\) adds general preference without explicit binding, and \(A2\) adds explicit binding.
        Deltas are paired by matched instance.
        Exp4 tests whether artifact action can persist without preference or binding, and whether preference or binding further changes artifact action.
        Because all three conditions are UC-side manipulations, this table does not by itself establish that the original UC--E gap persists after preference removal.
    \end{minipage}
\end{table}

\paragraph{Remaining preference-confound check.}
Exp4 shows that artifact action can persist without general preference, but it does not directly test whether the original \(UC-E\) difference persists after preference removal.
The minimal additional control is a matched external no-preference condition, \(E0\), with the same original task and artifact but no general preference and no explicit binding.
The key contrast would be \(A0_{UC}-E0\).
If \(E0 \ll A0_{UC}\), then the user-context effect persists after removing the general preference.
If \(E0 \approx A0_{UC}\), then the original \(UC-E\) difference should be reframed as depending substantially on preference, task relevance, or interaction structure.
